\documentclass[12pt, a4paper]{article}

% ============================================================
% Pakete (Stil identisch zu dokumentation.tex)
% ============================================================
\usepackage[T1]{fontenc}
\usepackage[utf8]{inputenc}
\usepackage[ngerman]{babel}
\usepackage{geometry}
\geometry{margin=2.5cm}
\setlength{\headheight}{15pt}

\usepackage{lmodern}
\usepackage{microtype}
\usepackage{parskip}
\usepackage{xcolor}
\usepackage{listings}
\usepackage{booktabs}
\usepackage{tabularx}
\usepackage{array}
\usepackage{hyperref}
\usepackage{fancyhdr}
\usepackage{titlesec}
\usepackage{enumitem}
\usepackage{tcolorbox}
\usepackage{amsmath}
\usepackage{amssymb}
\usepackage{float}
\usepackage{caption}
\usepackage{tikz}
\usetikzlibrary{shapes.geometric, arrows.meta, positioning, fit, backgrounds, calc}

% ============================================================
% Farben
% ============================================================
\definecolor{codebg}{RGB}{245,245,245}
\definecolor{codeframe}{RGB}{200,200,200}
\definecolor{keywordcolor}{RGB}{0,0,180}
\definecolor{commentcolor}{RGB}{100,130,80}
\definecolor{stringcolor}{RGB}{180,0,0}
\definecolor{accentblue}{RGB}{0,90,160}
\definecolor{accentgray}{RGB}{90,90,90}
\definecolor{boxbg}{RGB}{235,245,255}
\definecolor{boxframe}{RGB}{0,90,160}
\definecolor{tablegreen}{RGB}{225,245,230}
\definecolor{tablegreenframe}{RGB}{30,130,70}

% ============================================================
% Listings-Stil
% ============================================================
\lstset{
  literate=%
    {ä}{{\"{a}}}1 {ö}{{\"{o}}}1 {ü}{{\"{u}}}1
    {Ä}{{\"{A}}}1 {Ö}{{\"{O}}}1 {Ü}{{\"{U}}}1
    {ß}{{\ss{}}}1
    {→}{{$\rightarrow$}}1
    {←}{{$\leftarrow$}}1
    {…}{{...}}1
}

\lstdefinestyle{datastyle}{
  language=,
  backgroundcolor=\color{codebg},
  frame=single,
  rulecolor=\color{codeframe},
  basicstyle=\ttfamily\small,
  showstringspaces=false,
  breaklines=true,
  captionpos=b,
  xleftmargin=1.5em,
}

\lstdefinestyle{yamlstyle}{
  language=,
  backgroundcolor=\color{codebg},
  frame=single,
  rulecolor=\color{codeframe},
  basicstyle=\ttfamily\small,
  commentstyle=\color{commentcolor}\itshape,
  showstringspaces=false,
  breaklines=true,
  captionpos=b,
  xleftmargin=1.5em,
  morecomment=[l]{\#},
}

% ============================================================
% tcolorbox-Stile
% ============================================================
\tcbuselibrary{skins, breakable}

\newtcolorbox{infobox}[1][]{
  colback=boxbg,
  colframe=boxframe,
  fonttitle=\bfseries,
  title={#1},
  breakable,
}

\newtcolorbox{hinweisbox}[1][]{
  colback=tablegreen,
  colframe=tablegreenframe,
  fonttitle=\bfseries,
  title={#1},
  breakable,
}

% ============================================================
% Seitenlayout
% ============================================================
\pagestyle{fancy}
\fancyhf{}
\rhead{\textcolor{accentgray}{\small KI-Copilot für Verzahnungswissen}}
\lhead{\textcolor{accentgray}{\small CSV-Modul: Funktionsweise des Readers}}
\cfoot{\thepage}
\renewcommand{\headrulewidth}{0.4pt}

\hypersetup{
  colorlinks=true,
  linkcolor=accentblue,
  urlcolor=accentblue,
  pdftitle={CSV-Modul: Funktionsweise des Tabellen-Readers},
}

\titleformat{\section}{\Large\bfseries\color{accentblue}}{}{0em}{\thesection\quad}[\titlerule]
\titleformat{\subsection}{\large\bfseries\color{accentblue}}{}{0em}{\thesubsection\quad}
\titleformat{\subsubsection}{\normalsize\bfseries}{}{0em}{\thesubsubsection\quad}

% ============================================================
% Dokument
% ============================================================
\begin{document}

\begin{center}
  {\LARGE\bfseries\color{accentblue} CSV-Modul: Funktionsweise des Tabellen-Readers\par}
  \vspace{0.3cm}
  {\large\color{accentgray} Von der hochgeladenen Tabellen-Datei zur belegten Antwort\par}
  \vspace{0.6cm}
\end{center}

\begin{tcolorbox}[colback=boxbg, colframe=boxframe]
  Dieses Dokument erklärt, wie der KI-Copilot CSV- und Excel-Dateien verarbeitet:
  in welchen zwei Rollen eine Tabelle ins System gelangt, wie Zeilen für die
  Vektorsuche aufbereitet werden, wie die Suche Tabellen anders behandelt als
  PDF-Text und wie aus einer Zahnrad-CSV strukturierte Bauteildaten entstehen.
  Referenzierte Module: \texttt{tabular\_loader\_pandas.py},
  \texttt{chunker\_recursive.py}, \texttt{retriever\_hybrid.py},
  \texttt{csv\_gear\_mapper.py}.
\end{tcolorbox}

\tableofcontents
\newpage

% ============================================================
\section{Motivation: Warum Tabellen ein Sonderfall sind}
% ============================================================

Das RAG-System wurde ursprünglich für Fließtext (PDF-Normen, Fachliteratur)
entworfen. Tabellen brechen dabei gleich mehrere Grundannahmen:

\begin{enumerate}[itemsep=2pt]
  \item \textbf{Embedding-Lücke:} Das Embedding-Modell \texttt{bge-m3} ist auf
        natürliche Sprache optimiert. Nackte CSV-Zeilen wie
        \texttt{W-103,Ausgangswelle,1,42CrMo4} erzielen nur Ähnlichkeitswerte
        um $0{,}4$ -- und fallen damit unter den globalen Retrieval-Schwellenwert
        von $0{,}5$, bevor das LLM sie je zu sehen bekommt.
  \item \textbf{Kontextverlust beim Chunking:} Ein Satz-Chunker verschmilzt
        benachbarte Zeilen oder zerschneidet sie -- der Zeilenzusammenhang
        (\glqq eine Zeile = ein Datensatz\grqq) geht verloren.
  \item \textbf{Starre Treffergrenze:} Mit $k=5$ Treffern kann eine Frage wie
        \glqq Zeige mir \emph{alle} Wellen\grqq{} niemals vollständig beantwortet
        werden, sobald die Tabelle mehr als fünf relevante Zeilen enthält.
  \item \textbf{Semantische Verallgemeinerung:} Ohne strikte Spaltenbindung
        macht das LLM aus einer \glqq Ausgangswelle\grqq{} schnell eine
        generische \glqq Antriebswelle\grqq.
\end{enumerate}

Der Reader begegnet dem mit vier ineinandergreifenden Bausteinen:
\emph{Zeile$\rightarrow$Satz-Preprocessing} (Abschnitt~\ref{sec:preprocessing}),
\emph{tabellenbewusstem Chunking} (Abschnitt~\ref{sec:chunking}),
\emph{sortenspezifischen Schwellenwerten mit Hybrid-Suche}
(Abschnitt~\ref{sec:retrieval}) und einem \emph{Split-Search-Routing}
(Abschnitt~\ref{sec:routing}).

% ============================================================
\section{Zwei Rollen -- der Upload-Ort entscheidet}
% ============================================================

Eine CSV-Datei kann zwei grundverschiedene Dinge bedeuten: \emph{Wissen}
(z.\,B. eine Stückliste, die durchsuchbar sein soll) oder \emph{Bauteildaten}
(die Parameter genau eines Zahnrads, analog zu einer STEP-Analyse). Das System
entscheidet nicht per Heuristik, sondern über den Ort des Uploads:

\begin{figure}[H]
\centering
\begin{tikzpicture}[
  node distance=0.55cm and 1.0cm,
  box/.style={rectangle, rounded corners=2pt, draw=accentgray, fill=codebg,
              text width=4.6cm, align=center, font=\small, inner sep=6pt},
  chan/.style={rectangle, rounded corners=2pt, draw=boxframe, fill=boxbg,
               text width=4.6cm, align=center, font=\small, inner sep=6pt},
  arr/.style={-{Stealth}, thick, color=accentgray},
]
  \node[box] (lib) {\textbf{Dokumentbibliothek}\\ \glqq Dokumente hochladen\grqq};
  \node[box, right=2.2cm of lib] (comp) {\textbf{Composer / Sidebar}\\ Plus-Button bzw. \glqq Aktives Bauteil\grqq};

  \node[chan, below=of lib]  (know) {\texttt{POST /upload}\\ Kanal A: \textbf{Wissen}\\ Indexierung in Qdrant};
  \node[chan, below=of comp] (part) {\texttt{POST /cad/from-csv}\\ Kanal B: \textbf{Bauteildaten}\\ \texttt{cad\_metadata} für \texttt{/ask}};

  \draw[arr] (lib) -- (know);
  \draw[arr] (comp) -- (part);
\end{tikzpicture}
\caption{Rollentrennung: identische Datei, zwei Verarbeitungswege.}
\end{figure}

Beide Kanäle akzeptieren \texttt{.csv}, \texttt{.xlsx} und \texttt{.xls}.
PDF-Dateien laufen unabhängig vom Upload-Ort immer in die Wissensbasis.

% ============================================================
\section{Kanal A: CSV als Wissen}
% ============================================================

\subsection{Einlesen (\texttt{tabular\_loader\_pandas.py})}

Der \texttt{TabularLoader} implementiert das \texttt{DocumentLoader}-Protokoll
und liest die Datei robust ein:

\begin{itemize}[itemsep=2pt]
  \item \textbf{Encoding-Kaskade:} \texttt{utf-8} $\rightarrow$ \texttt{latin1}
        $\rightarrow$ \texttt{cp1252} $\rightarrow$ \texttt{iso-8859-1} --
        deckt auch Excel-Exporte mit Windows-Kodierung ab.
  \item \textbf{Separator-Erkennung:} \texttt{pandas} mit \texttt{sep=None}
        (Python-Engine) erkennt Komma, Semikolon und Tabulator automatisch.
  \item \textbf{Excel:} \texttt{.xlsx}/\texttt{.xls} werden über
        \texttt{openpyxl} gelesen.
  \item \textbf{Identität:} Der SHA-256-Hash der Datei wird zum
        \texttt{doc\_hash} -- Grundlage für Deduplikation und Löschung.
\end{itemize}

\subsection{Zeile$\rightarrow$Satz-Preprocessing}
\label{sec:preprocessing}

Kernidee: Jede Datenzeile wird beim Laden in eine natürlichsprachliche Aussage
umgewandelt, die \emph{gleichzeitig} die Spaltenstruktur explizit erhält.

\begin{lstlisting}[style=datastyle, caption={Vorher: Rohzeile aus \texttt{getriebe\_stueckliste.csv}}]
Bauteil_ID,Bezeichnung,Menge,Werkstoff,Modul_mm,Zaehnezahl,Oberflaechenhaerte
G-001,Stirnrad Antriebswelle,1,16MnCr5,3.0,24,58 HRC
\end{lstlisting}

\begin{lstlisting}[style=datastyle, caption={Nachher: Satzform, wie sie eingebettet und indexiert wird}]
Bauteil G-001 (Bauteil_ID: G-001): Bezeichnung = Stirnrad Antriebswelle; Menge = 1; Werkstoff = 16MnCr5; Modul_mm = 3.0; Zaehnezahl = 24; Oberflaechenhaerte = 58 HRC.
\end{lstlisting}

Die Konstruktionsregeln:

\begin{enumerate}[itemsep=2pt]
  \item \textbf{Subjekt aus der ersten Spalte:} Der Wert der ersten Spalte wird
        zum Satzanfang. Endet der Spaltenname auf ein ID-Suffix
        (\texttt{\_ID}, \texttt{-ID}, \dots), wird daraus das Subjektwort
        abgeleitet: \texttt{Bauteil\_ID} $\rightarrow$ \glqq Bauteil G-001\grqq.
  \item \textbf{Expliziter Header-Bezug:} Die Klammer
        \texttt{(Bauteil\_ID: G-001)} nennt den exakten Spaltennamen -- das
        ermöglicht dem LLM später strenges ID-Matching in der richtigen Spalte.
  \item \textbf{\glqq Spalte = Wert\grqq-Paare:} Alle weiteren gefüllten Spalten
        folgen als Paare mit Semikolon-Trennung. Die Spaltenzuordnung bleibt
        damit in jedem einzelnen Chunk vollständig erhalten.
  \item \textbf{Leerwerte entfallen:} \texttt{-}, \texttt{n/a}, \texttt{NaN},
        leere Zellen usw. werden übersprungen -- eine Welle ohne Modul erhält
        schlicht keine Modul-Aussage (statt eines irreführenden Platzhalters).
\end{enumerate}

\begin{infobox}[Warum das die Such-Scores hebt]
Embedding-Modelle bewerten Ähnlichkeit zwischen \emph{Sätzen}. Die Frage
\glqq Welchen Werkstoff hat W-103?\grqq{} und die Aussage
\glqq Bauteil W-103 (\dots): Werkstoff = 42CrMo4; \dots\grqq{} teilen sich
Vokabular \emph{und} Satzgestalt -- die Kosinus-Ähnlichkeit steigt deutlich
gegenüber der rohen Kommaliste. Zusätzlich enthält jede Zeile die Fachbegriffe
der Header, wodurch auch der lexikalische Sparse-Kanal
(Abschnitt~\ref{sec:retrieval}) exakt greift.
\end{infobox}

Das Ergebnis wird als \texttt{RawDocument} mit der Markierung
\texttt{doc\_kind = "table"} zurückgegeben -- dieses eine Feld steuert die
gesamte tabellenspezifische Sonderbehandlung stromabwärts.

\subsection{Tabellenbewusstes Chunking}
\label{sec:chunking}

Für Fließtext gilt weiterhin das satzbasierte Chunking nach Token-Limit
(bzw. der semantische Chunker). Trägt ein Dokument jedoch
\texttt{doc\_kind = "table"}, greifen \emph{beide} Chunker-Implementierungen
auf die gemeinsame Funktion \texttt{chunk\_table\_rows()} zurück:

\begin{itemize}[itemsep=2pt]
  \item \textbf{Eine Zeile = ein Chunk.} Kein Verschmelzen benachbarter
        Datensätze, kein Zerschneiden am Token-Limit, kein Overlap.
  \item \textbf{Keine Mindestgröße.} Der \texttt{min\_chunk\_tokens}-Filter
        (der kurze Zeilen verwerfen würde) ist für Tabellen deaktiviert.
  \item \textbf{Reihenfolge bleibt erhalten.} Die \texttt{position} zählt
        monoton -- beim späteren Voll-Laden erscheint die Tabelle in
        Originalreihenfolge.
\end{itemize}

Beim Indexieren wandert \texttt{doc\_kind} in die Metadaten jedes Chunks und
liegt damit als Payload-Feld in der Vektordatenbank Qdrant -- die Suche kann
Tabellenzeilen später von PDF-Text unterscheiden.

\subsection{Suche: Hybrid-Score und sortenspezifische Schwellenwerte}
\label{sec:retrieval}

Jede Frage wird mit \texttt{bge-m3} eingebettet und in Qdrant gesucht. Bei
aktivierter Hybrid-Suche kombiniert der Retriever den dichten (semantischen)
Score mit einem lexikalischen Sparse-Score (Bag-of-words, exakte
Begriffstreffer für IDs, Normnummern, Werkstoffcodes):

\begin{equation*}
  s(c) \;=\; w_d \cdot s_{\text{dense}}(c) \;+\; w_s \cdot s_{\text{sparse}}(c),
  \qquad w_d = 0{,}7,\; w_s = 0{,}3 .
\end{equation*}

Statt nur $k=5$ Kandidaten holt die Suche zunächst
$k \cdot \texttt{candidate\_multiplier} = 5 \cdot 8 = 40$ dichte Kandidaten und
sortiert sie nach dem Hybrid-Score neu.

Entscheidend für Tabellen ist der \textbf{sortenspezifische Schwellenwert}:
Tabellenzeilen erzielen systembedingt niedrigere Scores als Fließtext und
werden deshalb gegen eine eigene, niedrigere Schwelle geprüft:

\begin{equation*}
  \text{akzeptiert}(c) \;\Longleftrightarrow\;
  s(c) \;\geq\;
  \begin{cases}
    \tau_{\text{table}} = 0{,}35 & \text{falls } \texttt{doc\_kind}(c) = \texttt{table},\\[2pt]
    \tau_{\text{text}} = 0{,}50  & \text{sonst (PDF/Fließtext).}
  \end{cases}
\end{equation*}

Technisch läuft das zweistufig: Qdrant filtert serverseitig grob mit
$\min(\tau_{\text{text}}, \tau_{\text{table}})$ -- sonst wären Tabellenzeilen
verworfen, bevor der Retriever sie sieht --, danach wendet der Retriever die
sortenspezifische Schwelle auf den finalen Hybrid-Score an. Für PDF-Wissen
ändert sich dadurch nichts: Textchunks müssen unverändert $0{,}5$ erreichen.

\subsection{Split-Search-Routing: die \texttt{top\_k}-Grenze aufheben}
\label{sec:routing}

Nach Filterung entscheidet der Retriever \emph{deterministisch} (Regeln, kein
LLM-Aufruf), ob statt der besten $k$ Treffer eine ganze Tabelle als Kontext
geladen wird:

\begin{figure}[H]
\centering
\begin{tikzpicture}[
  node distance=0.5cm and 0.6cm,
  stage/.style={rectangle, rounded corners=2pt, draw=accentgray, fill=codebg,
               text width=5.6cm, align=center, font=\small, inner sep=5pt},
  dec/.style={rectangle, rounded corners=2pt, draw=boxframe, fill=boxbg,
              text width=5.6cm, align=center, font=\small, inner sep=5pt},
  res/.style={rectangle, rounded corners=2pt, draw=tablegreenframe, fill=tablegreen,
              text width=4.6cm, align=center, font=\small, inner sep=5pt},
  arr/.style={-{Stealth}, thick, color=accentgray},
  lab/.style={font=\scriptsize\itshape, color=accentgray},
]
  \node[stage] (search) {Hybrid-Suche: 40 Kandidaten,\\ sortenspezifisch gefiltert};
  \node[dec, below=of search] (kw) {\textbf{Route 1:} Enthält die Frage Listen-Signalwörter?\\ (\glqq alle\grqq, \glqq sämtliche\grqq, \glqq wie viele\grqq, \glqq Liste\grqq, \dots) \emph{und} gibt es einen Tabellen-Treffer?};
  \node[dec, below=of kw] (dom) {\textbf{Route 2 (ergebnisgetrieben):} Stammt die \emph{Mehrheit} der Top-$k$-Treffer aus \emph{derselben} Tabelle?};
  \node[dec, below=of dom] (size) {Hat diese Tabelle höchstens \texttt{auto\_table\_context\_rows} = 40 Zeilen?};
  \node[res, below=of size] (topk) {Normaler Pfad:\\ beste $k=5$ Treffer};
  \node[res, right=1.4cm of kw, yshift=-1.1cm] (full) {\textbf{Ganze Tabelle laden}\\ alle Zeilen in Originalreihenfolge\\ (Deckel: 150 Zeilen)};

  \draw[arr] (search) -- (kw);
  \draw[arr] (kw) -- node[lab, left] {nein} (dom);
  \draw[arr] (dom) -- node[lab, left] {ja} (size);
  \draw[arr] (size) -- node[lab, left] {nein} (topk);
  \draw[arr] (dom.west) -- node[lab, above=1pt] {nein} ++(-1.0,0) |- (topk.west);
  \draw[arr] (kw.east) to[out=0, in=140] node[lab, above right=-1pt] {ja} (full.north west);
  \draw[arr] (size.east) to[out=0, in=-90] node[lab, below right=-1pt] {ja} (full.south);
\end{tikzpicture}
\caption{Entscheidungsbaum des Split-Search-Routings (pro Frage, ohne LLM).}
\end{figure}

\textbf{Route 1 -- Signalwörter:} Ein regulärer Ausdruck erkennt
Aggregat-Formulierungen (u.\,a. \emph{alle, sämtliche, Liste, auflisten,
aufzählen, Übersicht, wie viele, Anzahl, gesamt, jede/jeder, vollständig,
komplett}; Wortgrenzen verhindern Fehltreffer wie \glqq Met\emph{alle}\grqq).
Trifft ein Signalwort zu und liegt unter den Kandidaten ein
Tabellen-Treffer, wird dessen Dokument komplett geladen
(max. \texttt{aggregate\_max\_chunks} = 150 Zeilen).

\textbf{Route 2 -- ergebnisgetrieben:} Formulierungen \emph{ohne} Signalwort
(\glqq \emph{Welche} Wellen gibt es?\grqq) erkennt keine Wortliste. Deshalb
wertet der Retriever zusätzlich das \emph{Suchergebnis selbst} aus: Stammt die
Mehrheit ($>50\,\%$) der Top-$k$-Treffer aus derselben Tabelle, \glqq landet\grqq{}
die Frage offensichtlich breit in dieser Tabelle -- viele Zeilen passen ähnlich
gut, und die starre Top-5-Auswahl würde willkürlich Zeilen verschlucken. In
diesem Fall wird die Tabelle ebenfalls komplett geladen -- aber nur, wenn sie
höchstens 40 Zeilen hat (Größenprüfung über ein Scroll-Sentinel). Große
Tabellen bleiben ohne explizites Listen-Signal beim Top-$k$-Pfad, damit
Einzelfakt-Fragen den LLM-Kontext nicht fluten.

Beim Voll-Laden erben alle Zeilen den Score des besten Tabellen-Treffers
(\emph{Seed}) -- für die Quellenanzeige zählt die Relevanz des Dokuments, nicht
die der Einzelzeile. Die Zeilen erscheinen in Originalreihenfolge.

\begin{hinweisbox}[Bewusster Nebeneffekt]
Besteht der Korpus (fast) nur aus einer kleinen Tabelle, bekommt auch eine
Einzelfakt-Frage die ganze Tabelle als Kontext -- die richtige Zeile ist dann
\emph{garantiert} enthalten, unabhängig davon, ob sie es in die Top-5 geschafft
hätte. Bei gemischtem Korpus mit dominanten PDF-Treffern bleibt der normale
Pfad erhalten.
\end{hinweisbox}

\subsection{Verhaltensregeln im LLM-Prompt}

Damit das Sprachmodell die aufbereiteten Zeilen korrekt nutzt, enthält der
Systemprompt tabellenspezifische Regeln:

\begin{itemize}[itemsep=2pt]
  \item \textbf{Strikte Spalten-Zuordnung:} Werte ausschließlich hinter dem
        exakt passenden Spaltennamen ablesen (\glqq Spalte = Wert\grqq-Format).
  \item \textbf{Kein Kategorien-Raten:} Bezeichnungen wortwörtlich übernehmen
        (\glqq Ausgangswelle\grqq{} wird nie zu \glqq Antriebswelle\grqq{}
        verallgemeinert).
  \item \textbf{Exaktes ID-Matching:} Eine angefragte ID (z.\,B. W-103) muss
        im ID-Feld eines Datensatzes stehen; existiert sie nicht, wird das
        bestätigt statt ähnliche IDs zu erfinden.
  \item \textbf{Mengen-Logik:} Fehlende Stückzahlen heißen
        \glqq nicht spezifiziert\grqq{} -- keine erfundenen Platzhalter.
  \item \textbf{Vollständige Auflistungen:} Bei Listen-/Zählfragen werden
        \emph{alle} zutreffenden Datensätze einzeln aufgezählt (kein
        \glqq usw.\grqq), mit Gesamtanzahl am Ende.
\end{itemize}

% ============================================================
\section{Kanal B: CSV als Bauteildaten}
% ============================================================

\subsection{Grundprinzip}

Der zweite Kanal behandelt eine CSV wie das Ergebnis einer CAD-Analyse:
\texttt{csv\_gear\_mapper.py} liest die Datei und erzeugt dieselbe
verschachtelte \texttt{GearParameters}-Struktur, die auch der CAD-Prozessor
für STEP-Dateien liefert. Das Ergebnis fließt unverändert als
\texttt{cad\_metadata} in \texttt{/ask} ein und wird im Inspector des Frontends
angezeigt.

\begin{hinweisbox}[Wichtig: Lesen, nicht Erkennen]
Anders als bei STEP wird hier \emph{keine Geometrie analysiert}. Der Mapper
liest ausschließlich, was in den Spalten steht -- er rechnet nichts nach und
rät nichts. Fehlt z.\,B. eine Typ-Spalte, bleibt \texttt{gear\_type} leer;
aus \glqq Modul 2{,}5 + 32 Zähne\grqq{} wird \emph{nicht} auf
\glqq Stirnrad\grqq{} geschlossen. Konvention: \textbf{eine CSV = ein Bauteil}
(bei Stücklisten wird die erste Zeile mit Verzahnungsdaten verwendet und eine
Warnung ausgegeben).
\end{hinweisbox}

\subsection{Spaltenerkennung}

Spaltennamen werden normalisiert (Kleinschreibung, Umlaute
\texttt{ä$\rightarrow$ae} usw., Einheiten-Klammern und Sonderzeichen entfernt:
\glqq Modul (mm)\grqq{} $\rightarrow$ \texttt{modul}) und gegen eine
Alias-Tabelle geprüft. Sie deckt deutsche und englische Namen sowie
DIN-Kurzzeichen ab:

\begin{table}[H]
\centering
\small
\begin{tabularx}{\textwidth}{@{}l X@{}}
\toprule
\textbf{Zielfeld (GearParameters)} & \textbf{akzeptierte Spaltennamen (Auswahl)} \\
\midrule
\texttt{gear\_type}            & Verzahnungstyp, Zahnradtyp, Bauart, Typ \\
\texttt{module\_mm}            & Modul, Modul\_mm, Module, Normalmodul, m \\
\texttt{num\_teeth}            & Zähnezahl, Zahnzahl, Teeth, z \\
\texttt{pressure\_angle\_deg}  & Eingriffswinkel, pressure\_angle, alpha \\
\texttt{helix\_angle\_deg}     & Schrägungswinkel, helix\_angle, beta \\
\texttt{pitch\_diameter\_mm}   & Teilkreis(durchmesser), d \\
\texttt{outer\_diameter\_mm}   & Kopfkreis, Außendurchmesser, da \\
\texttt{root\_diameter\_mm}    & Fußkreis, df \\
\texttt{face\_width\_mm}       & Zahnbreite, Breite, b \\
\texttt{material}              & Werkstoff, Material \\
\texttt{tolerance\_class}      & Härte, Oberflächenhärte, Toleranzklasse \\
\texttt{part\_name} / \texttt{part\_number} & Bezeichnung, Name / Bauteil\_ID, Teilenummer, Artikelnummer \\
\bottomrule
\end{tabularx}
\caption{Spalten-Alias-Zuordnung (erweiterbar in \texttt{csv\_gear\_mapper.py}).}
\end{table}

Mindestens eine der Signal-Spalten \emph{Modul}, \emph{Zähnezahl} oder
\emph{Verzahnungstyp} muss vorhanden sein -- sonst lehnt der Endpunkt die Datei
mit HTTP~400 und der Meldung \glqq keine erkennbaren Verzahnungs-Spalten\grqq{} ab.

\subsection{Wert-Parsing und Typ-Zuordnung}

\begin{itemize}[itemsep=2pt]
  \item \textbf{Zahlen:} Deutsches Dezimalkomma (\texttt{"2,5"}) und angehängte
        Einheiten (\texttt{"20°"}, \texttt{"25 mm"}) werden toleriert;
        Zähnezahlen werden zu Ganzzahlen gerundet.
  \item \textbf{Typbezeichnungen} werden auf das interne Enum abgebildet:
        Stirnrad\,$\rightarrow$\,\texttt{spur},
        Schrägverzahnung\,$\rightarrow$\,\texttt{helical},
        Kegelrad\,$\rightarrow$\,\texttt{bevel},
        Hohlrad/Innenverzahnung\,$\rightarrow$\,\texttt{internal},
        Schnecke\,$\rightarrow$\,\texttt{worm},
        Zahnstange\,$\rightarrow$\,\texttt{rack}.
        Unbekannte Begriffe werden unverändert durchgereicht.
  \item \textbf{Konfidenz:} Jeder Wert wird im Schema-2.0-Format
        \texttt{\{value, unit, confidence\}} ausgegeben, mit
        $\text{confidence} = 0{,}92$ -- das entspricht der Stufe
        \textsc{Direct} des CAD-Prozessors: \emph{explizit angegeben, nicht aus
        Geometrie geschätzt}.
\end{itemize}

\begin{lstlisting}[style=datastyle, caption={Auszug der Antwort von \texttt{POST /cad/from-csv}}]
{
  "schema_version": "2.0",
  "filename": "zahnrad.csv",
  "gear_type": {"value": "spur", "unit": "", "confidence": 0.92},
  "tooth_profile": {
    "module_mm":  {"value": 2.5, "unit": "mm", "confidence": 0.92},
    "num_teeth":  {"value": 32,  "unit": "",   "confidence": 0.92}
  },
  "material_context": {"material": "16MnCr5"},
  "metadata": {"part_number": "G-010", "created_by": "csv_gear_mapper"},
  "extraction_quality": {"warnings": []}
}
\end{lstlisting}

\subsection{Ablauf im GUI}

Nach dem Upload über den Plus-Button (Composer) oder \glqq Aktives Bauteil\grqq{}
(Sidebar) zeigt der Bauteil-Chip den Status (\emph{Wird analysiert} \dots{}
\emph{ready}), der Inspector rendert Modul, Zähnezahl, Werkstoff usw., und jede
Folgefrage sendet die Parameter als \texttt{cad\_metadata} mit. Antworten, die
sich darauf stützen, markiert das LLM mit \texttt{[CAD]}; Wissens-Treffer
tragen weiterhin \texttt{[Q1]}, \texttt{[Q2]}, \dots

% ============================================================
\section{Konfigurationsreferenz}
% ============================================================

Alle Stellschrauben des Readers liegen im Abschnitt \texttt{retriever} der
\texttt{config.yaml}:

\begin{lstlisting}[style=yamlstyle]
retriever:
  top_k: 5                     # Treffer im Normalpfad
  min_similarity: 0.5          # Schwelle für PDF-/Textchunks
  table_min_similarity: 0.35   # niedrigere Schwelle NUR für Tabellenzeilen
  use_hybrid: true             # Dense + lexikalischer Sparse-Kanal
  hybrid_dense_weight: 0.7
  hybrid_sparse_weight: 0.3
  candidate_multiplier: 8      # 40 Kandidaten vor dem Reranking
  aggregate_max_chunks: 150    # Deckel Route 1 (Signalwörter)
  auto_table_context_rows: 40  # Maximalgröße Route 2 (0 = Route 2 aus)
\end{lstlisting}

\begin{table}[H]
\centering
\small
\begin{tabularx}{\textwidth}{@{}l l X@{}}
\toprule
\textbf{Parameter} & \textbf{Default} & \textbf{Wirkung / wann anpassen} \\
\midrule
\texttt{table\_min\_similarity} & 0.35 & Tabellenzeilen fallen trotz Satzform durch $\rightarrow$ senken; zu viel Tabellen-Rauschen $\rightarrow$ anheben. \\
\texttt{aggregate\_max\_chunks} & 150 & Obergrenze für Voll-Laden bei Listenfragen; an Kontextfenster des LLM koppeln. \\
\texttt{auto\_table\_context\_rows} & 40 & Ab welcher Tabellengröße Route 2 nicht mehr automatisch voll lädt; 0 deaktiviert Route 2 komplett. \\
\texttt{candidate\_multiplier} & 8 & Breite der Vorsuche -- Basis der Mehrheitsentscheidung von Route 2. \\
\bottomrule
\end{tabularx}
\caption{Die vier tabellenrelevanten Stellschrauben.}
\end{table}

% ============================================================
\section{Grenzen und betriebliche Hinweise}
% ============================================================

\begin{itemize}[itemsep=2pt]
  \item \textbf{Reindex erforderlich:} Nach Aktivierung von Sparse-Vektoren
        bzw. Einführung von \texttt{doc\_kind} müssen bestehende Dokumente neu
        indexiert werden (\texttt{storage/qdrant} leeren, Dateien erneut
        hochladen) -- alte Qdrant-Punkte tragen die neuen Felder nicht.
  \item \textbf{Signalwortliste ist deutsch:} Route 1 erkennt deutsche
        Formulierungen; englische Fragen fallen auf Route 2 bzw. den
        Top-$k$-Pfad zurück.
  \item \textbf{Schwellenwerte sind Startwerte:} $\tau_{\text{table}} = 0{,}35$
        wurde konservativ gewählt; die tatsächlichen \texttt{bge-m3}-Scores der
        Satzform sollten im Live-Betrieb geprüft und der Wert ggf. nachjustiert
        werden.
  \item \textbf{Eine CSV = ein Bauteil (Kanal B):} Mehrzeilige Stücklisten
        gehören in die Wissensbasis; der Bauteil-Mapper nimmt sonst nur die
        erste Verzahnungszeile und warnt.
  \item \textbf{Kein Konversationskontext:} Wie im Gesamtsystem ist jede
        Anfrage isoliert -- die Tabelle wird pro Frage neu gesucht bzw.
        geladen.
\end{itemize}

\end{document}
